\section{Rationale}

Task-related EEG studies in ADHD have examined event-related changes in oscillatory activity across theta, alpha, and beta ranges, but the literature remains heterogeneous and does not support treating a simple spectral pattern as a diagnostic signature \parencite{michelini2022eventoscillations,slater2022eegsubtypes}. A broader review of EEG feature extraction in pediatric ADHD similarly shows substantial variation in paradigms and feature families, with spectral power and event-related potentials used for different executive-function questions \parencite{arpaia2022eegfeatures}.

ARC-ADHD creates an unusually explicit temporal structure for a measurement study: stable task performance, a controlled challenge, an overt regulation sequence, a defined recovery-monitoring state, and a behaviorally defined return to stable performance. The present paper asks whether wearable EEG dynamics track this sequence at the within-child event level.

Muse S Athena is a plausible measurement platform because the manufacturer currently describes real-time EEG sensing with seven EEG sensors and a 256~Hz acquisition rate for the device \parencite{interaxonAthena2026}. The manufacturer also provides a research pathway for academic users and states that its SDK/research tools can expose sensor data and manage acquisition \parencite{interaxonDevelopers2026}. Exact Athena-specific raw channel labels, montage, timestamp semantics, quality indicators, and licensed research access must nevertheless be confirmed before data collection.

The headband is not part of the core ARC-ADHD intervention. Its scientific value is conditional on measurement quality and incremental information. A null EEG result therefore leaves the behavioral challenge-regulation-recovery hypothesis conceptually intact.
